\documentclass[11pt]{article}
\usepackage[a4paper,margin=1in]{geometry}
\usepackage[T1]{fontenc}
\usepackage{lmodern,microtype}
\usepackage{amsmath,amssymb,amsthm,mathtools}
\usepackage{booktabs,enumitem}
\usepackage{xcolor}
\usepackage[numbers,sort&compress]{natbib}
\usepackage[colorlinks=true,linkcolor=blue!55!black,citecolor=blue!55!black]{hyperref}

\newtheorem{theorem}{Theorem}[section]
\newtheorem{proposition}[theorem]{Proposition}
\theoremstyle{remark}
\newtheorem{remark}[theorem]{Remark}
\newcommand{\norm}[1]{\left\lVert#1\right\rVert}

\title{Trace Powers versus Hilbert--Schmidt Mass in Nonnormal Cloud Complements}
\author{Liang Wang}
\date{July 2026}

\begin{document}
\maketitle

\begin{abstract}
RH-229 found that the whole-complement Frobenius estimate grows roughly like
$\sigma^{-1}$ and cannot prove a normal relative determinant family.  We
show that this failure is not a determinant obstruction.  The $N$-dimensional
nilpotent shift has $\norm{S_N}_2^2=N-1$, while
$\operatorname{tr}(S_N^n)=0$ for every $n\ge1$ and
$\det_2(I-zS_N)\equiv1$.  Nonnormal singular-value mass can therefore diverge
without contributing any spectral determinant mass.

For each RH-222 endpoint we compute the exact finite trace of the squared
scaled matrix and subtract the Perron, parity, and selected-cloud squares.
The maximum complement second-trace modulus is $0.12952$, whereas the
inherited Hilbert--Schmidt squared upper bound reaches $308.75$.  The largest
ratio is $2.36\times10^6$.  Hence the Frobenius route is too strong; the next
meaningful target is a uniform envelope for the cloud-extracted power traces.
\end{abstract}

\section{Two different notions of complement size}

The regularized determinant has logarithmic germ
\[
 \log\det_2(I-zB)=-\sum_{n\ge2}\frac{z^n}{n}\operatorname{tr}(B^n).
\]
An ideal-norm argument controls these traces through
$|\operatorname{tr}(B^n)|\le\norm{B^n}_1$, for example
$\norm{B^2}_1\le\norm B_2^2$ \citep{Simon2005}.  The inequality is sufficient
but can be extremely wasteful for nonnormal operators.

RH-229 measured precisely this sufficient budget and found it divergent
\citep{WangRH229}.  The present question is whether the divergence is spectral
or merely singular-value mass.

\section{Exact nilpotent separation}

\begin{theorem}[Divergent Hilbert--Schmidt mass with trivial determinant]
Let $S_N\in\mathbb C^{N\times N}$ be the shift with ones on the first
superdiagonal and zero elsewhere.  Then
\[
 \norm{S_N}_2^2=N-1,
 \qquad
 \operatorname{tr}(S_N^n)=0\quad(n\ge1),
 \qquad
 \det_2(I-zS_N)=1.
\]
In particular $\norm{S_N}_2\to\infty$ while the regularized determinant is
identically constant.
\end{theorem}

\begin{proof}
There are $N-1$ unit entries, proving the Hilbert--Schmidt identity.  The
matrix is strictly upper triangular, as is every positive power, so every
power trace vanishes.  All eigenvalues are zero, hence the canonical
regularized product equals one.
\end{proof}

This example is maximally nonnormal, but that is exactly its purpose: no
logical implication from a divergent Hilbert--Schmidt norm to determinant
divergence can hold without additional structure.

\begin{theorem}[Nontrivial determinant with an invisible HS sector]
Let $B$ be any fixed finite matrix and let $a_N\in\mathbb C$.  Define
\[
 A_N=B\oplus a_NS_N.
\]
Then, for every $n\ge1$,
\[
 \operatorname{tr}(A_N^n)=\operatorname{tr}(B^n),
 \qquad
 \det_2(I-zA_N)=\det_2(I-zB),
\]
whereas
\[
 \norm{A_N}_2^2=\norm B_2^2+|a_N|^2(N-1).
\]
Thus the Hilbert--Schmidt mass can diverge while an arbitrary fixed
nontrivial regularized determinant is preserved exactly.
\end{theorem}

\begin{proof}
Powers and traces split over direct sums.  The nilpotent block contributes
zero to every power trace and has regularized determinant one.  The
Hilbert--Schmidt identity is the Pythagorean formula for block direct sums.
\end{proof}

This extension shows that the separation is not tied to the constant
determinant one.  An arbitrarily complicated fixed spectral factor may
coexist with an increasingly large nonnormal singular-value sector that is
invisible to all power traces.

\section{Second-trace extraction}

Let $A_{\sigma,d}$ denote either the fine or Haar-coarse scaled matrix.  If
$p_\sigma$, $q_\sigma$, and $C_\sigma$ denote the Perron root, parity root,
and selected cloud, define
\[
 \tau_{2,\sigma}
 =\operatorname{tr}(A_{\sigma,d}^2)
  -p_\sigma^2-q_\sigma^2
  -\sum_{\lambda\in C_\sigma}\lambda^2.
\]
For a sparse real matrix,
\[
 \operatorname{tr}(A^2)=\sum_{i,j}A_{ij}A_{ji},
\]
so this coefficient can be computed without diagonalizing the unresolved
matrix.

\begin{table}[ht]
\centering
\small
\begin{tabular}{lr}
\toprule
Batch statistic & value\\
\midrule
Endpoint count & 32\\
Maximum $|\tau_{2,\sigma}|$ & $0.1295197$\\
Minimum $|\tau_{2,\sigma}|$ & $3.8778\times10^{-5}$\\
Maximum complement HS-squared upper & $308.7521$\\
Maximum HS-squared/$|\tau_2|$ ratio & $2.3619\times10^6$\\
\bottomrule
\end{tabular}
\caption{Spectral trace versus singular-value budget.}
\end{table}

The full scaled Frobenius mass grows strongly as the mesh resolves the narrow
kernel.  By contrast, $\operatorname{tr}(A^2)$ itself remains near $1.633$ at
the finest scales, and the selected cloud removes nearly all of the remaining
spectral second moment.

\section{Cancellation efficiency in the archived family}

It is useful to record the dimensionless ratio
\[
 \mathcal E_{2,\sigma}
 =\frac{\text{complement HS-squared upper}}
        {|\tau_{2,\sigma}|}.
\]
This is not a condition number and becomes infinite when the denominator
vanishes.  It measures only how much larger the cancellation-blind sufficient
budget is than the signed spectral coefficient that actually enters
$\log\det_2$ at order two.  Across the batch it ranges from ordinary
two-digit values at the coarsest endpoints to $2.36\times10^6$.  The sign of
$\tau_2$ also changes with noise and channel, confirming that phase and
nonnormal cancellation, rather than positive mass, govern this coefficient.

The sparse identity
$\operatorname{tr}(A^2)=\sum_{i,j}A_{ij}A_{ji}$ avoids an unresolved
eigendecomposition, but subtraction can still lose relative accuracy when
$\tau_2$ is tiny.  Accordingly the archive treats the values as frozen
double-precision observations.  A theorem would require interval control of
the full trace and the extracted cloud sum, or an analytic trace formula in
which their cancellation is performed before numerical evaluation.

\section{Why one trace coefficient is not enough}

The second coefficient is the first nontrivial term of the regularized
logarithm, so its smallness is necessary for convergence to a normalized
limit with vanishing quadratic coefficient.  It is not sufficient for a
normal determinant family.  Higher traces may remain large even when
$\tau_2=0$, and no bound on $\tau_n$ for $n\ge3$ follows from signed
cancellation at order two.

There are therefore two logically separate lessons:
\begin{enumerate}[leftmargin=2em]
\item failure of a Hilbert--Schmidt bound does not close the determinant
route;
\item success at one or finitely many trace orders does not close it either.
\end{enumerate}
The useful replacement is an order-uniform, noise-uniform trace envelope,
not merely a smaller norm surrogate.

\section{What the separation buys}

The nilpotent theorem and the finite audit jointly show that the route
\[
 \text{uniform HS complement}
 \Longrightarrow
 \text{normal relative determinant}
\]
is sufficient but not necessary.  It is therefore legitimate to replace the
failed norm gate by a direct trace gate.

This replacement has a cost.  One must control every $n\ge2$, not merely
$n=2$.  RH-236 computes orders through twelve and RH-240 formulates an
all-order geometric criterion.  Until such a bound is proved, no locally
uniform determinant family follows.

No claim about a Hilbert--Polya operator, zeta zeros, or the Riemann
Hypothesis is made.

\bibliographystyle{plainnat}
\bibliography{references}
\end{document}
